\section{Pregunta de investigación}

¿En qué medida la integración de un modelo matemático basado en ecuaciones diferenciales ordinarias, sensores IoT de bajo costo y algoritmos de aprendizaje automático permite desarrollar un sistema híbrido que no solo estime con precisión las emisiones de \ce{CO2} y \ce{CH4}, sino que constituya una herramienta de monitoreo dinámico para fundamentar la toma de decisiones operativas y generar métricas verificables del desempeño ambiental en la bioconversión de residuos agroindustriales con \textit{Hermetia illucens}?
